
\needspace{3\baselineskip}
\color{uniblue}{\subsection{Контроль привода РПН}}
\color{black}

\begin{enumerate}[label=\arabic{section}.\arabic{subsection}.\arabic{enumi}, labelsep=4pt, leftmargin=0pt, itemindent=57pt, itemsep=0pt, parsep=5pt] % если что можно крутить itemsep

%%%%%%%%%%%%%%%%%%%%%%%%%%%%%%%%%%%%%%%%%%%%%%%%%%%%%%%%%%%%%%%%%%%%% ОБЩИЕ СВЕДЕНИЯ %%%%%%%%%%%%%%%%%%%%%%%%%%%%%%%%%%%%%%%%%%%%%%%%%%%%%%%%%%%%%%

\item Общие сведения

\begin{enumerate}[label=\arabic{section}.\arabic{subsection}.\arabic{enumi}.\arabic*, labelsep=4pt, leftmargin=0em, itemindent=65pt, parsep=0pt]

\item Данный блок логики отвечает за:
\begin{itemize}
    \item формирование команд на привод (п. \ref{form});
    \item контроль привода в процессе переключения (п. \ref{kntr});
    \item формирование команд автоматического управления (п. \ref{avt});
    \item счетчик переключений и контроль выработанного ресурса (п. \ref{resurs}).
\end{itemize}

\item
В таблице \ref{kntr:tbl1} приведены параметры, необходимые для настройки блока логики контроля привода РПН.

\small
\begin{longtable}{|>{\centering\arraybackslash}m{5.3cm}|>{\centering\arraybackslash}m{3.3cm}|>{\centering\arraybackslash}m{4.2cm}|>{\centering\arraybackslash}m{1.8cm}|>{\centering\arraybackslash}m{1cm}|}
\caption{Параметры для настройки блока логики контроля привода РПН\hfill\vspace{-0.5\baselineskip}}\label{kntr:tbl1}\\ 
\hline
\rowcolor{gray!30}
Параметр (Параметр на ИЧМ) & Условное обозначение на схеме & Значение/ Диапазон & Единица измерения & Шаг \\ 
\hline
\endfirsthead
\caption*{\hspace{3pt}\emph{Продолжение таблицы \ref{kntr:tbl1}\hfill\vspace{-0.5\baselineskip}}} \\ % сделано по ГОСТ 2.105 п.6.8.7
\hline
\rowcolor{gray!30}
Параметр (Параметр на ИЧМ) & Условное обозначение на схеме & Значение/ Диапазон & Единица измерения & Шаг \\ 
%\hline
\endhead
%\hline
\endfoot
\hline
\endlastfoot
%==+t1*ATCC1|ARNT_TTCC
\centering Источник сигнала положения РПН (Датчик\_РПН)  & \centering Датчик\_РПН & \centering  & \centering -- & \centering \arraybackslash -- \\
\hline
\centering Первоначальное значение счетчика положений РПН (Нач\_полож\_РПН)  & \centering Нач\_полож\_РПН & \centering 1...40 & \centering -- & \centering \arraybackslash 1 \\
\hline
\centering Номер крайнего нижнего положения РПН (Мин\_положение\_РПН)  & \centering Мин\_положение\_РПН & \centering 1...40 & \centering -- & \centering \arraybackslash 1 \\
\hline
\centering Номер крайнего верхнего положения РПН (Макс\_положение\_РПН)  & \centering Макс\_положение\_РПН & \centering 1...40 & \centering -- & \centering \arraybackslash 1 \\
\hline
\centering Направление счета положений (Направление\_счета)  & \centering Направление\_счета & \centering  & \centering -- & \centering \arraybackslash -- \\
\hline
\centering Сигнал токовой петли, соответствующий крайнему нижнему положению РПН (Iмин\_ТП\_РПН)  & \centering Iмин\_ТП\_РПН & \centering 0...4000 & \centering мкА & \centering \arraybackslash 1 \\
\hline
\centering Сигнал токовой петли, соответствующий крайнему верхнему положению РПН (Iмакс\_ТП\_ РПН)  & \centering Iмакс\_ТП\_ РПН & \centering 5000...20000 & \centering мкА & \centering \arraybackslash 1 \\
\hline
\centering Положение привода, далее которого блокируются автоматические команды прибавить (Верхн\_разреш\_полож)  & \centering Верхн\_разреш\_полож & \centering 1...40 & \centering -- & \centering \arraybackslash 1 \\
\hline
\centering Положение привода, далее которого блокируются автоматические команды убавить (Нижн\_разреш\_полож)  & \centering Нижн\_разреш\_полож & \centering 1...40 & \centering -- & \centering \arraybackslash 1 \\
\hline
\centering Промежуточное 1 положение РПН (Промежут\_1\_полож)  & \centering Промежут\_1\_полож & \centering 0...40 & \centering -- & \centering \arraybackslash 1 \\
\hline
\centering Промежуточное 2 положение РПН (Промежут\_2\_полож)  & \centering Промежут\_2\_полож & \centering 0...40 & \centering -- & \centering \arraybackslash 1 \\
\hline
\centering Промежуточное 3 положение РПН (Промежут\_3\_полож)  & \centering Промежут\_3\_полож & \centering 0...40 & \centering -- & \centering \arraybackslash 1 \\
\hline
\centering Промежуточное 4 положение РПН (Промежут\_4\_полож)  & \centering Промежут\_4\_полож & \centering 0...40 & \centering -- & \centering \arraybackslash 1 \\
\hline
\centering Промежуточное 5 положение РПН (Промежут\_5\_полож)  & \centering Промежут\_5\_полож & \centering 0...40 & \centering -- & \centering \arraybackslash 1 \\
\hline
\centering Центр зоны нечувствительности (Uпод)  & \centering Uпод & \centering 0,85...1,45 & \centering о.е. & \centering \arraybackslash 0,01 \\
\hline
\centering Ширина зоны нечувствительности (dUпод)  & \centering dUпод & \centering 0,01...0,20 & \centering о.е. & \centering \arraybackslash 0,01 \\
\hline
\centering Уровень напряжения, при превышении которого автоматические команды на снижение формируются с ускоренной выдержкой времени T3 (U>>сраб)  & \centering U>>сраб & \centering 1,00...2,00 & \centering о.е. & \centering \arraybackslash 0,01 \\
\hline
\centering Величина тока перегрузки РПН (Iрпн>блк)  & \centering Iрпн>блк & \centering 1,00...2,00 & \centering о.е. & \centering \arraybackslash 0,01 \\
\hline
\centering Блокировка от контролируемой секции (БлкКонтрСекц)  & \centering БлкКонтрСекц & \centering Не предусмотрено \\ Предусмотрено & \centering -- & \centering \arraybackslash -- \\
\hline
\centering Величина тока перегрузки для запрета авторегулирования на повышение (I>блк)  & \centering I>блк & \centering 1,000...2,000 & \centering о.е. & \centering \arraybackslash 0,001 \\
\hline
\centering Величина напряжения НП для запрета авторегулирования на повышение (3U0блк)  & \centering 3U0блк & \centering 0,00...2,00 & \centering о.е. & \centering \arraybackslash 0,01 \\
\hline
\centering Величина напряжения ОП для запрета авторегулирования (U2блк)  & \centering U2блк & \centering 0,00...2,00 & \centering о.е. & \centering \arraybackslash 0,01 \\
\hline
\centering Величина максимального напряжения для запрета авторегулирования на повышение (U>>блк)  & \centering U>>блк & \centering 1,00...2,00 & \centering о.е. & \centering \arraybackslash 0,01 \\
\hline
\centering Величина минимального напряжения для запрета авторегулирования (U<<блк)  & \centering U<<блк & \centering 0,00...1,00 & \centering о.е. & \centering \arraybackslash 0,01 \\
\hline
\centering Время задержки срабатывания на выдачу первой команды в автоматическом режиме (T1)  & \centering T1 & \centering 1000...200000 & \centering мс & \centering \arraybackslash 1 \\
\hline
\centering Время задержки срабатывания на выдачу второй и последующих команд в автоматическом режиме (T2)  & \centering T2 & \centering 100...200000 & \centering мс & \centering \arraybackslash 1 \\
\hline
\centering Время задержки выдачи ускоренной команды убавить при значительном повышении напряжения в автоматическом режиме (T3)  & \centering T3 & \centering 100...10000 & \centering мс & \centering \arraybackslash 1 \\
\hline
\centering Выдержка времени снятия команды управления после появления сигнала переключения привода (Твозвр)  & \centering Твозвр & \centering 0...10000 & \centering мс & \centering \arraybackslash 1 \\
\hline
\centering Максимальное время реакции привода РПН на команду управления (Тпуск)  & \centering Тпуск & \centering 100...10000 & \centering мс & \centering \arraybackslash 1 \\
\hline
\centering Максимальное время одного переключения привода РПН (Тперекл)  & \centering Тперекл & \centering 100...100000 & \centering мс & \centering \arraybackslash 1 \\
\hline
\centering Режим команды на отключение питания ПМ при самопроизвольном переключении (Режим\_откл\_ПМ)  & \centering Режим\_откл\_ПМ & \centering  & \centering -- & \centering \arraybackslash -- \\
\hline
\centering Ресурс привода РПН по количеству операций переключения (Ресурс\_переключений)  & \centering Ресурс\_переключений & \centering 1...1000000 & \centering -- & \centering \arraybackslash 1 \\
\hline
\centering Начальное значение счетчика переключений привода РПН (Нач\_знач\_счетчика)  & \centering Нач\_знач\_счетчика & \centering 1...1000000 & \centering -- & \centering \arraybackslash 1 \\
\hline
\centering Метод параллельного регулирования (Метод\_парал\_рег)  & \centering Метод\_парал\_рег & \centering  & \centering -- & \centering \arraybackslash -- \\
\hline
\centering Контроль рассогласования приводов РПН трансформаторов в режиме параллельного регулирования (Контроль\_рассогл)  & \centering Контроль\_рассогл & \centering Внешний сигнал \\ Сигналы положений приводов РПН & \centering -- & \centering \arraybackslash -- \\
\hline
\centering Источник данных о положении РПН параллельного трансформатора (Датчик\_парал\_РПН)  & \centering Датчик\_парал\_РПН & \centering Токовая петля 4(0)-20мА \\ GOOSE \\ BCD-код & \centering -- & \centering \arraybackslash -- \\
\hline
\centering Максимально допустимый разбег приводов РПН при параллельном регулировании трансформаторов (Макс\_разбег\_РПН)  & \centering Макс\_разбег\_РПН & \centering 0...40 & \centering -- & \centering \arraybackslash 1 \\
\hline
\centering Время переключения РПН трансформаторов при параллельном регулировании (Тпарал\_перекл)  & \centering Тпарал\_перекл & \centering 100...100000 & \centering мс & \centering \arraybackslash 1 \\
\hline
\centering Номер крайнего нижнего положения РПН параллельного трансформатора (Мин\_полож\_пар\_РПН)  & \centering Мин\_полож\_пар\_РПН & \centering 1...40 & \centering -- & \centering \arraybackslash 1 \\
\hline
\centering Номер крайнего верхнего положения РПН параллельного трансформатора (Макс\_полож\_пар\_РПН)  & \centering Макс\_полож\_пар\_РПН & \centering 1...40 & \centering -- & \centering \arraybackslash 1 \\
\hline
\centering Сигнал токовой петли, соответствующий крайнему нижнему положению РПН параллельного трансформатора (Iмин\_ТП\_парал\_РПН)  & \centering Iмин\_ТП\_парал\_РПН & \centering 0...4000 & \centering мкА & \centering \arraybackslash 1 \\
\hline
\centering Сигнал токовой петли, соответствующий крайнему верхнему положению РПН параллельного трансформатора (Iмакс\_ТП\_парал\_РПН)  & \centering Iмакс\_ТП\_парал\_РПН & \centering 5000...20000 & \centering мкА & \centering \arraybackslash 1 \\
\hline
\centering Выдержка времени срабатывания сигнализации (Тсигн)  & \centering Тсигн & \centering 1...100000 & \centering мс & \centering \arraybackslash 1 \\
\hline
\centering Контроль положения блока испытательного №1 (Контроль\_БИ\_1)  & \centering Контроль\_БИ\_1 & \centering Не предусмотрено \\ Предусмотрено & \centering -- & \centering \arraybackslash -- \\
\hline
\centering Контроль положения блока испытательного №2 (Контроль\_БИ\_2)  & \centering Контроль\_БИ\_2 & \centering Не предусмотрено \\ Предусмотрено & \centering -- & \centering \arraybackslash -- \\
\hline
\centering Контроль положения блока испытательного №3 (Контроль\_БИ\_3)  & \centering Контроль\_БИ\_3 & \centering Не предусмотрено \\ Предусмотрено & \centering -- & \centering \arraybackslash -- \\
\hline
\centering Контроль положения блока испытательного №4 (Контроль\_БИ\_4)  & \centering Контроль\_БИ\_4 & \centering Не предусмотрено \\ Предусмотрено & \centering -- & \centering \arraybackslash -- \\
\hline
\centering Контроль положения блока испытательного №5 (Контроль\_БИ\_5)  & \centering Контроль\_БИ\_5 & \centering Не предусмотрено \\ Предусмотрено & \centering -- & \centering \arraybackslash -- \\
\hline
\centering Контроль положения блока испытательного №6 (Контроль\_БИ\_6)  & \centering Контроль\_БИ\_6 & \centering Не предусмотрено \\ Предусмотрено & \centering -- & \centering \arraybackslash -- \\
\hline
\centering Контроль положения блока испытательного №7 (Контроль\_БИ\_7)  & \centering Контроль\_БИ\_7 & \centering Не предусмотрено \\ Предусмотрено & \centering -- & \centering \arraybackslash -- \\
\hline
\centering Контроль положения двери (Контроль\_двери)  & \centering Контроль\_двери & \centering Не предусмотрено \\ Предусмотрено & \centering -- & \centering \arraybackslash -- \\
\hline
\centering Контроль оперативного тока технологической сигнализации (Контроль\_опер\_тока)  & \centering Контроль\_опер\_тока & \centering Не предусмотрено \\ Предусмотрено & \centering -- & \centering \arraybackslash -- \\
%===t1
\end{longtable}
\normalsize

\end{enumerate}

%%%%%%%%%%%%%%%%%%%%%%%%%%%%%%%%%%%%%%%%%%%%%%%%%%%%%%%%%%%%%%%%%%%%% Формирование команд на привод %%%%%%%%%%%%%%%%%%%%%%%%%%%%%%%%%%%%%%%%%%%%%%%%%%
\item\label{form} Формирование команд на привод

\begin{enumerate}[label=\arabic{section}.\arabic{subsection}.\arabic{enumi}.\arabic*, labelsep=4pt, leftmargin=0em, itemindent=65pt, parsep=0pt]

\item В зависимости от выбранного режима работы устройство может формировать три вида команд управления на привод:
\begin{itemize}
    \item оперативные;
    \item автоматические;
    \item от устройства регулирования трансформатора, включенного на параллельную работу.
\end{itemize}

\item\label{vydacha}
Оперативные команды управления имеют взаимную блокировку команд <<Прибавить>> и <<Убавить>>. Таймер <<t3>> (см. рисунок \ref{kntr:img1}) ограничивает длительность команды управления, в импульсном режиме сбрасывает команду после появления сигнала <<Идет переключение>>. В непрерывном режиме управления уставка данного таймера должна быть достаточной, чтобы привод гарантировано начал переключение.

\end{enumerate}

\vspace{2mm}
\begin{figure}[H]
\centering
\includegraphics[width=1\textwidth,height=1\textheight,keepaspectratio]{\fbpath/501. Контроль РПН/_latex/img/1.pdf}
\captionof{figure}{Функциональная схема логики выдачи команд управления на привод}\label{kntr:img1}
\end{figure}

%%%%%%%%%%%%%%%%%%%%%%%%%%%%%%%%%%%%%%%%%%%%%%%%%%%%%%%%%%%%%%%%%%%%% Контроль привода в процессе переключения %%%%%%%%%%%%%%%%%%%%%%%%%%%%%%%%%%%%%%%%%%%%%%%%%%

\item\label{kntr} Контроль привода в процессе переключения

\begin{enumerate}[label=\arabic{section}.\arabic{subsection}.\arabic{enumi}.\arabic*, labelsep=4pt, leftmargin=0em, itemindent=65pt, parsep=0pt]

\item Функциональная схема логики контроля привода приведена на рисунке \ref{kntr:img2}.

\item Кроме таймера <<t3>> в логике устройства предусмотрен ряд таймеров, используемых для управления приводом и для его диагностики в процессе переключения в случае импульсного режима выдачи команд:

\begin{itemize}
\item таймер <<t1>> (время реакции привода на команду управления) -- при отсутствии переключения в течение выдержки времени таймера <<t1>> формируется сигнал <<ПМ не пошел>>;
\item таймер <<t2>> (максимальная продолжительность переключения) -- при удержании сигнала <<Идет переключение>> в течение выдержки времени таймера <<t2>> устройство формирует сигнал <<ПМ застрял>>.
\end{itemize}

\item При отсутствии команд управления и появлении сигнала <<Идет переключение>> устройство формирует сигнал <<ПМ побежал>>. При возникновении данной ситуации, устройство ожидает завершение переключения и формирует сигнал <<Отключить автомат РПН>>.

\item
Программный переключатель SGF2 определяет время воздействия на дистанционный расцепитель автоматического выключателя привода:
\begin{itemize}
\item непрерывном;
\item в течение 1 с.
\end{itemize}

\item Выбор режима выдачи управляющих команд (Импульсный/Непрерывный) осуществляется программным переключателем SGF1.

\vspace{2mm}
\begin{figure}[H]
\centering
\includegraphics[width=1\textwidth,height=1\textheight,keepaspectratio]{\fbpath/501. Контроль РПН/_latex/img/2.pdf}
\captionof{figure}{Функциональная схема логики контроля привода}\label{kntr:img2}
\end{figure}

\end{enumerate}

%%%%%%%%%%%%%%%%%%%%%%%%%%%%%%%%%%%%%%%%%%%%%%%%%%%%%%%%%%%%%%%%%%%%% Формирование команд автоматического управления %%%%%%%%%%%%%%%%%%%%%%%%%%%%%%%%%%%%%%%%%%%%%%%%%%

\item\label{avt} Формирование команд автоматического управления

\begin{enumerate}[label=\arabic{section}.\arabic{subsection}.\arabic{enumi}.\arabic*, labelsep=4pt, leftmargin=0em, itemindent=65pt, parsep=0pt]

\item
В режиме автоматического регулирования устройство непрерывно отслеживает уровень напряжения на регулируемой секции (см. раздел \ref{avtreg}). При выходе напряжения за пределы <<Зоны нечувствительности>> устройство с определенными выдержками времени формирует сигналы управления приводом (см. рисунок \ref{kntr:img3}). 
\item
Используются следующие выдержки времени:
\begin{itemize}
    \item задержка на выдачу первой автоматической команды (уставка таймера <<T1>>) <<Прибавить>> или <<Убавить>> при выходе значения напряжения регулируемой секции за пределы зоны нечувствительности;
    \item задержка на выдачу повторной автоматической команды (уставка таймера <<T2>>) <<Прибавить>> или <<Убавить>>; вторая и все последующие команды выдаются с данной задержкой до возврата напряжения в пределы зоны нечувствительности;
    \item задержка на выдачу команды «Убавить» (уставка таймера <<T3>>) при значительном увеличении напряжения. При перенапряжении на регулируемой секции первая и все последующие автоматические команды <<Убавить>> выдаются с данной выдержкой времени.  
\end{itemize}

\vspace{2mm}
\begin{figure}[H]
\centering
\includegraphics[width=1\textwidth,height=1\textheight,keepaspectratio]{\fbpath/501. Контроль РПН/_latex/img/3.pdf}
\captionof{figure}{Функциональная схема логики формирования автоматических команд управления}\label{kntr:img3}
\end{figure}

\end{enumerate}

%%%%%%%%%%%%%%%%%%%%%%%%%%%%%%%%%%%%%%%%%%%%%%%%%%%%%%%%%%%%%%%%%%%%% Счетчик переключений и контроль выработанного ресурса %%%%%%%%%%%%%%%%%%%%%%%%%%%%%%%%%%%%%%%%%%%%%%%%%%

\item\label{resurs} Счетчик переключений и контроль выработанного ресурса

\begin{enumerate}[label=\arabic{section}.\arabic{subsection}.\arabic{enumi}.\arabic*, labelsep=4pt, leftmargin=0em, itemindent=65pt, parsep=0pt]

\item
В устройстве предусмотрен контроль выработанного ресурса. При каждом переключении значение счетчика переключений сравнивается с уставкой <<Максимальное количество операций переключения>>. При превышении уставки формируется сигнал превышения максимального количества переключений. Имеется возможность сбросить значение счетчика переключений через виртуальный ключ. Также имеется возможность задать уставкой начальное значение счетчика.

\item
Алгоритм работы счетчика переключений и контроля выработанного ресурса выполнен в соответствии с рисунком \ref{kntr:img4}.

\vspace{2mm}
\begin{figure}[H]
\centering
\includegraphics[width=1\textwidth,height=1\textheight,keepaspectratio]{\fbpath/501. Контроль РПН/_latex/img/4.pdf}
\captionof{figure}{Функциональная схема логики формирования автоматических команд управления}\label{kntr:img4}
\end{figure}

\end{enumerate}


\end{enumerate}